% Modelo físico-matemático del rastreador de cambios NISAR para Ecuador (estilo cyberpunk).
% Compilar desde la raíz del proyecto:
%   python docs/generar_datos_latex.py      (regenera datos_generados.tex)
%   cd docs && pdflatex modelo_nisar_ecuador.tex && pdflatex modelo_nisar_ecuador.tex
% Versión para imprimir (fondo blanco): \usepackage[impresion]{estilo_cyberpunk}
\documentclass[11pt,a4paper]{article}
\usepackage[margin=2.4cm,headsep=0.8cm]{geometry}
\usepackage[utf8]{inputenc}
\usepackage[T1]{fontenc}
\usepackage{lmodern}
\usepackage[spanish,es-noshorthands,es-tabla,es-nolists]{babel}
\usepackage{amsmath,amssymb}
\usepackage{booktabs}
\usepackage{graphicx}
\usepackage{microtype}
\usepackage{pgfplots}
\usepgfplotslibrary{groupplots}
\pgfplotsset{compat=1.18}
\usepackage{estilo_cyberpunk}
\usepackage{hyperref}

\newcommand{\dB}{\,\mathrm{dB}}
\newcommand{\vect}[1]{\mathbf{#1}}
\newcommand{\codigo}[1]{\cpcodigo{#1}}

% Archivo generado por docs/generar_datos_latex.py; no editar a mano.

\newcommand{\TablaClimaSitios}{%
\begin{table}[htbp]\centering\scriptsize
\caption{Climatología NASA POWER 2001--2025 interpolada a los sitios de Ecuador y temporadas resultantes. $P$ y $ET_0$ son totales anuales (mm); $z_c$ es la elevación media de la celda MERRA-2 (m), no la del sitio.}\label{tab:clima-sitios}
\setlength{\tabcolsep}{2.5pt}
\begin{tabular}{@{}lllrrrlp{2.5cm}p{2.5cm}@{}}\toprule
Sitio & Provincia & Región & $P$ & $ET_0$ & $z_c$ & Régimen & Meses lluviosos & Meses secos ($<60$ mm) \\ \midrule
Humedal La Segua & Manabí & Costa & 947 & 1585 & 129 & unimodal & ene, feb, mar, abr & jun, jul, ago, sep, oct, nov, dic \\
Bosques andinos de Loja & Loja & Sierra & 662 & 1446 & 1818 & unimodal & ene, feb, mar, abr & may, jun, jul, ago, sep, oct, nov, dic \\
Pedernales & Manabí & Costa & 1166 & 1400 & 136 & unimodal & ene, feb, mar, abr, may & jun, jul, ago, sep, oct, nov \\
Glaciares del volcán Antisana & Napo & Amazonía & 1423 & 976 & 3121 & unimodal & ene, feb, mar, abr, may, nov, dic & ago \\
Arrozales de Daule & Guayas & Costa & 801 & 1972 & 100 & unimodal & ene, feb, mar, abr & may, jun, jul, ago, sep, oct, nov, dic \\
Guayaquil & Guayas & Costa & 751 & 1837 & 237 & unimodal & ene, feb, mar, abr & may, jun, jul, ago, sep, oct, nov, dic \\
\bottomrule\end{tabular}\end{table}}

\newcommand{\FiguraClimaSitios}{%
\begin{figure}[htbp]\centering
\begin{tikzpicture}
\begin{groupplot}[group style={group size=3 by 2, horizontal sep=1.1cm, vertical sep=1.4cm},
  width=0.36\textwidth, height=4.2cm, ymin=0, xmin=0.5, xmax=12.5, xtick={1,...,12},
  xticklabels={E,F,M,A,M,J,J,A,S,O,N,D}, tick label style={font=\scriptsize},
  title style={font=\small}, ylabel style={font=\scriptsize}, ymajorgrids]
\nextgroupplot[title={Humedal La Segua}, ylabel={mm/mes}]
\addplot[ybar, bar width=4.2pt, fill=lluviosa, draw=none] coordinates {(1,162.6) (2,215.1) (3,216.2) (4,153.2)};
\addplot[ybar, bar width=4.2pt, fill=transicion, draw=none] coordinates {(5,74.2)};
\addplot[ybar, bar width=4.2pt, fill=seca, draw=none] coordinates {(6,27.2) (7,10.0) (8,5.7) (9,7.8) (10,12.9) (11,14.9) (12,47.6)};
\addplot[thick, dashed, cplinea, mark=*, mark size=0.9pt] coordinates {(1,129.5) (2,108.5) (3,119.1) (4,117.8) (5,119.2) (6,117.1) (7,132.2) (8,148.6) (9,152.1) (10,154.9) (11,145.0) (12,141.1)};
\addplot[densely dotted, cpumbral, domain=0.5:12.5, samples=2] {60};
\nextgroupplot[title={Bosques andinos de Loja}]
\addplot[ybar, bar width=4.2pt, fill=lluviosa, draw=none] coordinates {(1,66.2) (2,99.5) (3,103.7) (4,70.8)};
\addplot[ybar, bar width=4.2pt, fill=seca, draw=none] coordinates {(5,44.8) (6,38.6) (7,31.2) (8,23.0) (9,22.9) (10,47.4) (11,53.5) (12,59.8)};
\addplot[thick, dashed, cplinea, mark=*, mark size=0.9pt] coordinates {(1,126.0) (2,110.6) (3,119.1) (4,112.5) (5,110.6) (6,100.7) (7,108.7) (8,124.4) (9,134.2) (10,139.2) (11,132.3) (12,128.1)};
\addplot[densely dotted, cpumbral, domain=0.5:12.5, samples=2] {60};
\nextgroupplot[title={Pedernales}]
\addplot[ybar, bar width=4.2pt, fill=lluviosa, draw=none] coordinates {(1,194.3) (2,222.8) (3,234.4) (4,185.0) (5,112.6)};
\addplot[ybar, bar width=4.2pt, fill=transicion, draw=none] coordinates {(12,69.3)};
\addplot[ybar, bar width=4.2pt, fill=seca, draw=none] coordinates {(6,50.4) (7,19.5) (8,11.3) (9,17.3) (10,24.3) (11,24.4)};
\addplot[thick, dashed, cplinea, mark=*, mark size=0.9pt] coordinates {(1,112.8) (2,96.6) (3,107.1) (4,104.5) (5,103.7) (6,100.9) (7,115.1) (8,131.4) (9,134.9) (10,138.1) (11,129.8) (12,124.7)};
\addplot[densely dotted, cpumbral, domain=0.5:12.5, samples=2] {60};
\nextgroupplot[title={Glaciares del volcán Antisana}, ylabel={mm/mes}]
\addplot[ybar, bar width=4.2pt, fill=lluviosa, draw=none] coordinates {(1,132.9) (2,137.8) (3,167.2) (4,178.3) (5,142.9) (11,126.6) (12,123.4)};
\addplot[ybar, bar width=4.2pt, fill=transicion, draw=none] coordinates {(6,98.8) (7,76.5) (9,66.2) (10,114.5)};
\addplot[ybar, bar width=4.2pt, fill=seca, draw=none] coordinates {(8,58.2)};
\addplot[thick, dashed, cplinea, mark=*, mark size=0.9pt] coordinates {(1,87.5) (2,79.6) (3,88.1) (4,81.1) (5,77.5) (6,67.3) (7,70.9) (8,78.1) (9,83.7) (10,90.4) (11,85.9) (12,85.4)};
\addplot[densely dotted, cpumbral, domain=0.5:12.5, samples=2] {60};
\nextgroupplot[title={Arrozales de Daule}]
\addplot[ybar, bar width=4.2pt, fill=lluviosa, draw=none] coordinates {(1,128.6) (2,216.6) (3,214.9) (4,131.4)};
\addplot[ybar, bar width=4.2pt, fill=seca, draw=none] coordinates {(5,42.3) (6,12.0) (7,4.3) (8,3.4) (9,3.0) (10,7.2) (11,9.5) (12,27.4)};
\addplot[thick, dashed, cplinea, mark=*, mark size=0.9pt] coordinates {(1,163.4) (2,132.4) (3,141.2) (4,142.8) (5,151.3) (6,152.3) (7,169.4) (8,186.4) (9,189.7) (10,190.3) (11,177.5) (12,175.3)};
\addplot[densely dotted, cpumbral, domain=0.5:12.5, samples=2] {60};
\nextgroupplot[title={Guayaquil}]
\addplot[ybar, bar width=4.2pt, fill=lluviosa, draw=none] coordinates {(1,114.6) (2,195.0) (3,200.2) (4,118.8)};
\addplot[ybar, bar width=4.2pt, fill=seca, draw=none] coordinates {(5,39.6) (6,13.7) (7,6.9) (8,5.7) (9,5.5) (10,11.5) (11,13.7) (12,26.1)};
\addplot[thick, dashed, cplinea, mark=*, mark size=0.9pt] coordinates {(1,153.0) (2,125.3) (3,133.2) (4,134.9) (5,142.7) (6,142.9) (7,157.4) (8,172.5) (9,175.5) (10,175.1) (11,163.0) (12,161.9)};
\addplot[densely dotted, cpumbral, domain=0.5:12.5, samples=2] {60};
\end{groupplot}
\end{tikzpicture}
\caption{Precipitación mensual climatológica (barras: cian lluviosa, gris transición, magenta seca), $ET_0$ de Hargreaves--Samani (línea discontinua) y umbral de mes seco de 60~mm (punteada) en los seis sitios del catálogo.}\label{fig:clima-sitios}
\end{figure}}

\newcommand{\FiguraForzamiento}{%
\begin{figure}[htbp]\centering
\begin{tikzpicture}
\begin{groupplot}[group style={group size=1 by 2, vertical sep=0.9cm}, width=0.92\textwidth,
  height=4.3cm, xmin=0, xmax=365, xtick={0,31,59,90,120,151,181,212,243,273,304,334},
  xticklabels={ene,feb,mar,abr,may,jun,jul,ago,sep,oct,nov,dic}, ymajorgrids,
  tick label style={font=\scriptsize}, ylabel style={font=\small}, legend style={font=\scriptsize, at={(0.99,0.97)}, anchor=north east}]
\nextgroupplot[ylabel={\%}, ymin=0, ymax=105, xticklabels={}]
\addplot[thick, color=lluviosa] coordinates {(0,33.8) (3,38.2) (6,42.8) (9,47.9) (12,53.2) (15,58.8) (18,64.6) (21,70.4) (24,76.3) (27,82.2) (30,88.3) (33,94.3) (36,100.0) (39,100.0) (42,100.0) (45,100.0) (48,100.0) (51,100.0) (54,100.0) (57,100.0) (60,100.0) (63,100.0) (66,100.0) (69,100.0) (72,100.0) (75,100.0) (78,100.0) (81,100.0) (84,100.0) (87,100.0) (90,100.0) (93,100.0) (96,100.0) (99,100.0) (102,100.0) (105,100.0) (108,100.0) (111,100.0) (114,100.0) (117,100.0) (120,99.8) (123,99.2) (126,98.1) (129,96.6) (132,94.7) (135,92.4) (138,90.0) (141,87.4) (144,84.7) (147,82.0) (150,79.1) (153,76.2) (156,73.2) (159,70.1) (162,67.0) (165,63.8) (168,60.6) (171,57.6) (174,54.6) (177,51.7) (180,48.9) (183,46.1) (186,43.4) (189,40.8) (192,38.3) (195,35.8) (198,33.5) (201,31.2) (204,29.2) (207,27.2) (210,25.3) (213,23.6) (216,21.9) (219,20.4) (222,19.0) (225,17.6) (228,16.3) (231,15.2) (234,14.1) (237,13.2) (240,12.4) (243,11.6) (246,11.0) (249,10.4) (252,9.8) (255,9.4) (258,8.9) (261,8.6) (264,8.3) (267,8.1) (270,7.9) (273,7.8) (276,7.7) (279,7.7) (282,7.7) (285,7.7) (288,7.7) (291,7.8) (294,7.9) (297,8.0) (300,8.1) (303,8.2) (306,8.3) (309,8.4) (312,8.6) (315,8.7) (318,8.8) (321,9.1) (324,9.5) (327,10.1) (330,10.8) (333,11.7) (336,12.7) (339,13.8) (342,15.0) (345,16.3) (348,17.7) (351,19.4) (354,21.6) (357,24.3) (360,27.5) (363,31.2)};
\addlegendentry{humedad del suelo $s$ (porcentaje del agua útil)}
\addplot[thick, color=seca] coordinates {(0,0.0) (3,0.0) (6,0.0) (9,0.0) (12,0.0) (15,0.0) (18,0.0) (21,0.0) (24,0.0) (27,0.0) (30,0.0) (33,0.0) (36,0.9) (39,12.9) (42,24.8) (45,36.5) (48,47.2) (51,56.5) (54,64.7) (57,71.8) (60,78.0) (63,83.3) (66,87.8) (69,91.7) (72,94.9) (75,97.5) (78,99.1) (81,99.9) (84,99.9) (87,99.2) (90,97.8) (93,95.8) (96,93.3) (99,90.3) (102,86.8) (105,83.0) (108,78.6) (111,73.7) (114,68.2) (117,62.3) (120,56.3) (123,50.9) (126,46.0) (129,41.5) (132,37.5) (135,33.9) (138,30.6) (141,27.6) (144,25.0) (147,22.5) (150,20.4) (153,18.4) (156,16.6) (159,15.0) (162,13.6) (165,12.2) (168,11.1) (171,10.0) (174,9.0) (177,8.2) (180,7.4) (183,6.7) (186,6.0) (189,5.4) (192,4.9) (195,4.4) (198,4.0) (201,3.6) (204,3.3) (207,2.9) (210,2.7) (213,2.4) (216,2.2) (219,2.0) (222,1.8) (225,1.6) (228,1.4) (231,1.3) (234,1.2) (237,1.1) (240,1.0) (243,0.9) (246,0.8) (249,0.7) (252,0.6) (255,0.6) (258,0.5) (261,0.5) (264,0.4) (267,0.4) (270,0.3) (273,0.3) (276,0.3) (279,0.3) (282,0.2) (285,0.2) (288,0.2) (291,0.2) (294,0.2) (297,0.1) (300,0.1) (303,0.1) (306,0.1) (309,0.1) (312,0.1) (315,0.1) (318,0.1) (321,0.1) (324,0.1) (327,0.1) (330,0.0) (333,0.0) (336,0.0) (339,0.0) (342,0.0) (345,0.0) (348,0.0) (351,0.0) (354,0.0) (357,0.0) (360,0.0) (363,0.0)};
\addlegendentry{almacenamiento del humedal $S$ (porcentaje del máximo anual)}
\nextgroupplot[ylabel={$\Delta\sigma^0_{suelo}$ (dB)}]
\addplot[thick, cpamarillo] coordinates {(0,-0.63) (3,-0.47) (6,-0.29) (9,-0.11) (12,0.09) (15,0.31) (18,0.52) (21,0.74) (24,0.96) (27,1.18) (30,1.41) (33,1.64) (36,1.85) (39,1.85) (42,1.85) (45,1.85) (48,1.85) (51,1.85) (54,1.85) (57,1.85) (60,1.85) (63,1.85) (66,1.85) (69,1.85) (72,1.85) (75,1.85) (78,1.85) (81,1.85) (84,1.85) (87,1.85) (90,1.85) (93,1.85) (96,1.85) (99,1.85) (102,1.85) (105,1.85) (108,1.85) (111,1.85) (114,1.85) (117,1.85) (120,1.84) (123,1.82) (126,1.78) (129,1.72) (132,1.65) (135,1.57) (138,1.47) (141,1.38) (144,1.28) (147,1.17) (150,1.07) (153,0.96) (156,0.84) (159,0.73) (162,0.61) (165,0.49) (168,0.37) (171,0.26) (174,0.15) (177,0.04) (180,-0.07) (183,-0.17) (186,-0.27) (189,-0.37) (192,-0.46) (195,-0.56) (198,-0.65) (201,-0.73) (204,-0.81) (207,-0.88) (210,-0.95) (213,-1.02) (216,-1.08) (219,-1.14) (222,-1.19) (225,-1.24) (228,-1.29) (231,-1.33) (234,-1.37) (237,-1.40) (240,-1.44) (243,-1.46) (246,-1.49) (249,-1.51) (252,-1.53) (255,-1.55) (258,-1.56) (261,-1.58) (264,-1.59) (267,-1.60) (270,-1.60) (273,-1.61) (276,-1.61) (279,-1.61) (282,-1.61) (285,-1.61) (288,-1.61) (291,-1.61) (294,-1.60) (297,-1.60) (300,-1.60) (303,-1.59) (306,-1.59) (309,-1.58) (312,-1.58) (315,-1.57) (318,-1.57) (321,-1.56) (324,-1.54) (327,-1.52) (330,-1.50) (333,-1.46) (336,-1.43) (339,-1.38) (342,-1.34) (345,-1.29) (348,-1.24) (351,-1.17) (354,-1.09) (357,-0.99) (360,-0.87) (363,-0.73)};
\end{groupplot}
\end{tikzpicture}
\caption{Ciclo anual del forzamiento estacional en el humedal La Segua (Manabí): humedad del suelo del balance hídrico, almacenamiento del humedal (embalse lineal, $\tau=30$ días) y cambio resultante de la retrodispersión del suelo, ecuación~\eqref{eq:dsigma}.}\label{fig:forzamiento}
\end{figure}}

\newcommand{\TablaHitos}{%
\begin{table}[htbp]\centering\small
\caption{Hitos astronómicos según la declinación de Spencer (día juliano $J$ y fecha de 2026) y, para cada sitio, pasos del Sol por el cenit ($\delta=\varphi$) y duración del día en los solsticios (h).}\label{tab:hitos}
\begin{tabular}{@{}lrr@{}}\toprule Hito & $J$ & Fecha \\ \midrule
Equinoccio de marzo & 80{,}16 & 21/03 \\
Solsticio de junio & 173{,}01 & 22/06 \\
Equinoccio de septiembre & 266{,}63 & 24/09 \\
Solsticio de diciembre & 356{,}29 & 22/12 \\
\bottomrule\end{tabular}\par\bigskip
\begin{tabular}{@{}lrlrr@{}}\toprule Sitio & $\varphi$ ($^\circ$) & Pasos cenitales & $N_{jun}$ & $N_{dic}$ \\ \midrule
Humedal La Segua & $-$0{,}70 & 19/03 y 25/09 & 11{,}96 & 12{,}04 \\
Bosques andinos de Loja & $-$4{,}05 & 11/03 y 04/10 & 11{,}77 & 12{,}23 \\
Pedernales & $-$0{,}07 & 21/03 y 24/09 & 12{,}00 & 12{,}00 \\
Glaciares del volcán Antisana & $-$0{,}47 & 20/03 y 25/09 & 11{,}97 & 12{,}03 \\
Arrozales de Daule & $-$1{,}86 & 16/03 y 28/09 & 11{,}89 & 12{,}11 \\
Guayaquil & $-$2{,}19 & 16/03 y 29/09 & 11{,}87 & 12{,}13 \\
\bottomrule\end{tabular}\end{table}}

\newcommand{\TablaVerificacion}{%
\begin{table}[htbp]\centering\footnotesize\setlength{\tabcolsep}{4pt}
\caption{Verificación: indicadores de cada escenario de Ecuador sin forzamiento estacional (fase fija heredada) y con forzamiento climatológico. Semilla 7, 12 \emph{looks}, órbita ascendente.}\label{tab:verificacion}
\begin{tabular}{@{}llll@{}}\toprule Tipo & Sitio & Fase fija & Forzamiento estacional \\ \midrule
Pérdida de humedales & Humedal La Segua & máx. 31/03, mín. 14/11 & máx. 19/03, mín. 26/11 \\
Incendios forestales & Bosques andinos de Loja & amplitud HV 0{,}1 dB & amplitud HV 0{,}1 dB \\
Terremotos & Pedernales & $\sigma_{atm}$ constante, 1{,}80 cm & $\sigma_{atm}$ 1{,}69--1{,}89 cm \\
Movimiento de glaciares & Glaciares del volcán Antisana & retroceso 15 m & retroceso 15 m \\
Actividad agrícola & Arrozales de Daule & exactitud 98,9 \% & exactitud 98,9 \% \\
Subsidencia del terreno & Guayaquil & RMSE 1,98 cm/año & RMSE 1,93 cm/año \\
\bottomrule\end{tabular}\end{table}}

\newcommand{\TablaInvariancia}{%
\begin{table}[htbp]\centering\small
\caption{Prueba de invariancia: fecha del máximo del pulso de inundación en La Segua (sin tendencia de secado) al desplazar la fecha de referencia. Con fase fija el pico se desplaza con la fecha; con forzamiento permanece en la época lluviosa.}\label{tab:invariancia}
\begin{tabular}{@{}lrll@{}}\toprule Fecha de referencia & $\Delta$ (días) & Pico con fase fija & Pico con forzamiento \\ \midrule
01/05/2026 & +0 & 31/03 & 19/03 \\
30/07/2026 & +90 & 29/06 & 25/03 \\
28/10/2026 & +180 & 27/09 & 26/03 \\
\bottomrule\end{tabular}
\par\medskip
\begin{tabular}{@{}lrl@{}}\toprule Fecha de referencia (Daule) & $\Delta$ (días) & Meses de cosecha de arroz \\ \midrule
15/05/2026 & +0 & abr, may, sep, oct \\
12/09/2026 & +120 & abr, may, sep, oct \\
\bottomrule\end{tabular}\end{table}}

\newcommand{\TablaParticion}{%
\begin{table}[htbp]\centering\small
\caption{Resultado de la partición de demostración: 25 bloques de 8~km, buffer de 1500~m y semilla 20261011. Superficie efectiva tras el recorte por buffer.}\label{tab:particion}
\begin{tabular}{@{}lrrrr@{}}\toprule Conjunto & Bloques & Superficie (km$^2$) & \% superficie & Recortado (km$^2$) \\ \midrule
Entrenamiento & 15 & 708{,}91 & 52{,}6 & 251{,}09 \\
Validación & 5 & 320{,}00 & 23{,}7 & 0{,}00 \\
Prueba & 5 & 320{,}00 & 23{,}7 & 0{,}00 \\
Excluido por buffer & 0 & 0{,}00 & 0{,}0 & 0{,}00 \\
\bottomrule\end{tabular}
\par\smallskip{\scriptsize Huella SHA-256: \texttt{c61528b146153d03d755f9ba95832bfd}\allowbreak\texttt{557374fce2bf3af97593705eb1b86d39}}
\end{table}}
\newcommand{\BufferDemo}{1500}
\newcommand{\BloqueDemo}{8000}

\cptitulocorto{Modelo físico-matemático · señal SAR · temporadas}

\begin{document}
\cpportada{> NISAR // BANDA L // ECUADOR\_}
  {Modelo físico-matemático del rastreador de cambios NISAR para Ecuador}
  {Señal SAR, fenómenos superficiales y forzamiento estacional}
  {Documentación técnica del proyecto \emph{modelo clima}}
  {\today}

\begin{cpcaja}[title=RESUMEN]
Se documenta el modelo que sustenta la aplicación \codigo{app.py}: un simulador físico de las
observaciones de radar de apertura sintética (SAR) en banda L de la misión NISAR (NASA--ISRO) y los
algoritmos de detección de cambios que se aplican tanto a escenas simuladas como a productos reales
GCOV y GUNW. El dominio se restringe a Ecuador (continente y Galápagos). Se describen la geometría de
observación, el modelo estadístico de la señal (\emph{speckle}, nube de agua, ruido de fase y pantalla
atmosférica), seis modelos directos de fenómenos superficiales (humedales, incendios, terremotos,
subsidencia, glaciares y agricultura) y sus detectores. La contribución principal de esta versión es un
\emph{forzamiento estacional} explícito: para cada punto y fecha se determinan la región natural, la
estación astronómica y la temporada climática a partir de la climatología NASA POWER 2001--2025, y un
balance hídrico diario del suelo acopla esa temporada con la retrodispersión, el pulso de inundación, el
retardo troposférico y el calendario agrícola. Una prueba de invariancia demuestra que, a diferencia de
la formulación anterior, la estacionalidad queda ligada al calendario y no a la fecha del evento.
\end{cpcaja}

\tableofcontents
\clearpage

%=====================================================================================================
\section{Introducción y alcance}
%=====================================================================================================

NISAR observa la superficie terrestre con un radar en banda L ($f_L = 1{,}257$~GHz) y otro en banda S
($f_S = 3{,}2$~GHz) con un ciclo de repetición de 12~días. La aplicación permite estudiar seis tipos de
cambio en lugares de Ecuador; para cada uno combina:
\begin{enumerate}
  \item un \emph{modelo físico directo} que describe el estado de la superficie en cualquier instante
        $t$ (resolución de minutos);
  \item un \emph{modelo de observación} que transforma ese estado en lo que mediría NISAR en cada
        adquisición, con sus fuentes de error dominantes;
  \item un \emph{detector de cambios} idéntico al que se aplica a los productos reales;
  \item un \emph{forzamiento estacional} (sección~\ref{sec:estaciones}) que impone el ciclo anual del
        agua propio del lugar.
\end{enumerate}
Las escenas simuladas no son observaciones reales: reproducen la física y los errores de los
productos de nivel~2 para que los algoritmos puedan evaluarse contra una \emph{verdad} conocida.

En lo que sigue, $t$ denota días fraccionarios desde la fecha de referencia del escenario, $t_a = t/365{,}25$
el tiempo en años, $(X, Y)$ las coordenadas métricas locales (este, norte) de la escena y $G_\ell$ un
campo gaussiano estandarizado con longitud de correlación de $\ell$ píxeles (ecuación~\ref{eq:campo}).

%=====================================================================================================
\section{Dominio geográfico y regiones naturales}
%=====================================================================================================

El dominio es el territorio de Ecuador, representado por las 24 provincias de geoBoundaries
\cite{runfola2020} simplificadas a una tolerancia de $0{,}003^\circ$ ($\approx 300$~m). Cada provincia se
asigna a una región natural según la agrupación usual del INEC:
\begin{itemize}
  \item \textbf{Costa:} El Oro, Esmeraldas, Guayas, Los Ríos, Manabí, Santa Elena y Santo Domingo de los
        Tsáchilas;
  \item \textbf{Sierra:} Azuay, Bolívar, Cañar, Carchi, Chimborazo, Cotopaxi, Imbabura, Loja, Pichincha y
        Tungurahua;
  \item \textbf{Amazonía:} Morona Santiago, Napo, Orellana, Pastaza, Sucumbíos y Zamora Chinchipe;
  \item \textbf{Insular:} Galápagos.
\end{itemize}
Un punto $(\varphi,\lambda)$ pertenece a Ecuador si está contenido en alguna provincia o si su distancia
al contorno simplificado es menor que $0{,}02^\circ$ ($\approx 2$~km), tolerancia que absorbe la
simplificación de la línea de costa. La aplicación rechaza los clics fuera de este dominio y restringe el
mapa a la caja $[-92{,}3^\circ, -75{,}1^\circ]\times[-5{,}1^\circ, 1{,}8^\circ]$ (longitud $\times$
latitud) con un margen de $2^\circ$. La región natural es administrativa: un punto de alta montaña en
una provincia amazónica (p.~ej. el Antisana, en Napo) se clasifica como Amazonía, aunque su clima se
toma siempre de la climatología del propio punto.

%=====================================================================================================
\section{Geometría de observación NISAR}
%=====================================================================================================

Las longitudes de onda son $\lambda = c/f$, con $c = 299\,792\,458$~m\,s$^{-1}$:
$\lambda_L = 23{,}85$~cm y $\lambda_S = 9{,}37$~cm. La órbita es heliosíncrona con inclinación
$i = 98{,}4^\circ$ \cite{nisar2018}. El rumbo de la traza a la latitud $\varphi$ se aproxima por
\begin{equation}
  \alpha = \arcsin\!\left(\frac{\cos i}{\max(\cos\varphi,\, 0{,}05)}\right),\qquad
  h = \begin{cases} \alpha & \text{paso ascendente},\\ 180^\circ - \alpha & \text{paso descendente},\end{cases}
\end{equation}
y, como NISAR mira a la izquierda, el vector unitario suelo$\to$satélite en coordenadas (este, norte,
arriba) para un ángulo de incidencia $\theta$ es
\begin{equation}
  \hat{\vect e}_{LOS} = (\sin\theta\cos h,\; -\sin\theta\sin h,\; \cos\theta).
\end{equation}
El desplazamiento observado en la línea de vista es $d_{LOS} = \vect u\cdot\hat{\vect e}_{LOS}$. En
Ecuador ($\varphi\approx 0$), $|\alpha|\approx 8{,}4^\circ$: la componente norte del movimiento casi no se
proyecta en la línea de vista.

%=====================================================================================================
\section{Modelo de la señal radar}
%=====================================================================================================

\subsection{Campos aleatorios}
Las heterogeneidades del terreno se generan con ruido blanco $w$ filtrado en el dominio de Fourier y
estandarizado a media 0 y desviación 1:
\begin{equation}\label{eq:campo}
  G_\ell = \mathcal{N}\!\left\{\mathcal F^{-1}\!\left[\mathcal F[w]\, e^{-2\pi^2\ell^2|\vect k|^2}\right]\right\},
\end{equation}
donde $\mathcal N\{\cdot\}$ estandariza y $\vect k$ es la frecuencia espacial en ciclos por píxel.

\subsection{Speckle}
La intensidad multilook de $L$ vistas equivalentes sigue una distribución gamma \cite{goodman1976}:
\begin{equation}
  I = \sigma^0\,\xi,\qquad \xi\sim\Gamma\!\left(L,\, 1/L\right),\qquad \mathrm{E}[\xi]=1,\quad \mathrm{Var}[\xi]=1/L.
\end{equation}
Los detectores aplican un promedio móvil $3\times3$ (multilook espacial) que ignora los valores no
definidos.

\subsection{Modelo de nube de agua}
La retrodispersión de una cubierta vegetal con biomasa $B$ (Mg\,ha$^{-1}$) se modela con la nube de agua
de Attema y Ulaby \cite{attema1978}:
\begin{equation}\label{eq:wcm}
  \sigma^0 = A\cos\theta\,(1-\tau^2) + \sigma^0_{s}\,\tau^2,\qquad \tau^2 = \exp\!\left(-\frac{2bB}{\cos\theta}\right),
\end{equation}
con $\theta = 40^\circ$ y los parámetros de la tabla~\ref{tab:wcm}. Su inversión, usada para estimar la
biomasa perdida en incendios, es
\begin{equation}
  B = -\frac{\cos\theta}{2b}\ln\frac{A\cos\theta-\sigma^0}{A\cos\theta-\sigma^0_s},
\end{equation}
con $\sigma^0$ acotado a $[1{,}02\,\sigma^0_s,\; 0{,}985\,A\cos\theta]$ y $B\le 300$~Mg\,ha$^{-1}$. La saturación
de HV en banda L cerca de $100$--$150$~Mg\,ha$^{-1}$ hace que la biomasa perdida en bosques densos sea una
cota inferior.

\begin{table}[htbp]\centering\small
\caption{Parámetros del modelo de nube de agua (valores lineales).}\label{tab:wcm}
\begin{tabular}{@{}lrrr@{}}\toprule Cubierta y canal & $A$ & $\sigma^0_s$ & $b$ \\ \midrule
Bosque, HV & 0{,}075 & 0{,}004 & 0{,}006 \\
Bosque, HH & 0{,}20 & 0{,}03 & 0{,}006 \\
Cultivo, HV & 0{,}035 & 0{,}0025 & 0{,}10 \\
Cultivo, HH & 0{,}08 & 0{,}025 & 0{,}06 \\ \bottomrule
\end{tabular}
\end{table}

\subsection{Fase interferométrica}
La fase de un interferograma entre dos fechas es
\begin{equation}
  \Delta\phi = -\frac{4\pi}{\lambda}\,d_{LOS} + \phi_{atm} + \phi_{ruido},
\end{equation}
de modo que una franja ($2\pi$) equivale a $\lambda/2 = 11{,}9$~cm. La desviación del ruido de fase para
coherencia $\gamma$ y $L$ vistas viene dada por la cota de Cramér--Rao \cite{rodriguez1992}:
\begin{equation}
  \sigma_\phi = \frac{\sqrt{1-\gamma^2}}{\gamma\sqrt{2L}},\qquad \sigma_d = \frac{\lambda}{4\pi}\,\sigma_\phi ,
\end{equation}
con $\gamma$ acotada a $[0{,}05;\,0{,}999]$ y $4L$ vistas en el interferograma multilook.

\subsection{Pantalla atmosférica}
El retardo troposférico turbulento se modela como un campo con espectro de potencia de Kolmogórov
$P(k)\propto k^{-\beta}$, $\beta = 8/3$ \cite{williams1998,hanssen2001}, estandarizado y escalado por
$\sigma_{atm}$ y por el factor estacional $F_{pw}(t)$ de la ecuación~\eqref{eq:fpw}:
\begin{equation}
  a_k = \sigma_{atm}\,F_{pw}(t_k)\;\mathcal N\!\left\{\mathcal F^{-1}\!\left[\mathcal F[w_k]\,|\vect k|^{-\beta/2}\right]\right\}.
\end{equation}
Las pantallas son independientes entre adquisiciones; en la pila se observa $a_k - a_0$.

%=====================================================================================================
\section{Modelos directos de los fenómenos}
%=====================================================================================================

\subsection{Humedales}
La llanura se describe con un río meandriforme $x_r(Y) = A_m\sin(2\pi Y/\Lambda)$, la distancia al cauce
$D = |X - x_r|$ y una topografía estandarizada
\begin{equation}
  z = \mathcal N\!\left\{0{,}55\,G_{14} + 0{,}25\,G_4 + 1{,}4\left(1 - e^{-(D/3000)^2}\right)\right\}.
\end{equation}
El bosque ribereño ocupa los píxeles con $G_5 + 1{,}3\,e^{-(D/1200)^2} > 0{,}7$. El nivel del agua es
\begin{equation}\label{eq:nivel}
  h(t) = h_0 + a_e\,p(t) - r_s\,\frac{\max(t,0)}{365{,}25},
\end{equation}
donde $h_0$ es el cuantil de $z$ correspondiente a la fracción inundada inicial, $a_e$ la amplitud del
pulso estacional, $r_s$ la tasa de secado y $p(t)\in[-1,1]$ el pulso de inundación. Con forzamiento
estacional, $p(t) = \tilde S(t)$ es la anomalía normalizada del almacenamiento del humedal
(ecuación~\ref{eq:embalse}); sin él, $p(t)=\cos[2\pi(t-t_p)/365{,}25]$. Con la profundidad $h - z$ se
asignan cinco clases (agua abierta, herbácea inundada si la lámina es $\le 0{,}25$, bosque inundado, bosque
y herbácea), con retrodispersiones HH $= (-22, -7, -3, -8, -12)\dB$ y HV $= (-28, -17, -11, -12{,}5, -19)\dB$,
una textura $0{,}8\,G_{1,5}$~dB y una rugosidad del agua por viento $U(0,4)$~dB en HH (la mitad en HV).

\subsection{Incendios forestales}
La biomasa inicial es $B_0 = \mathrm{clip}(\mu_B + s_B G_8;\, 20;\, 450)$, con claros ($G_4 > 1{,}4$) de
$U(3,15)$~Mg\,ha$^{-1}$. El frente es una elipse que nace en el punto de ignición $(x_0,y_0)$ y crece
empujada por el viento de dirección $\psi$. En coordenadas normalizadas a favor del viento y
transversales,
\begin{equation}
  U = \frac{(X-x_0)\sin\psi + (Y-y_0)\cos\psi}{a},\qquad V = \frac{(X-x_0)\cos\psi - (Y-y_0)\sin\psi}{0{,}55\,a},
\end{equation}
la elipse de crecimiento $g$ tiene semiejes $(g a,\, 0{,}55\,g a)$ y centro $0{,}7\,g a$ a sotavento, por
lo que un píxel arde cuando $(U-0{,}7g)^2 + V^2 \le g^2$, es decir, cuando $g$ supera
\begin{equation}
  g_{\min} = \frac{-1{,}4\,U + \sqrt{1{,}96\,U^2 + 2{,}04\,(U^2+V^2)}}{1{,}02} \;(+\,0{,}12\,G_6).
\end{equation}
Con $g(t) = 1 - e^{-(t-t_i)/\tau}$ y $\tau = 2$~días, el instante de quema es
$t_q = t_i - \tau\ln(1-g_{\min})$ para $g_{\min}<1$ y $B_0>25$. La biomasa tras la quema es
$B = B_0\,[1 - \varsigma + \varsigma\,(1 - e^{-(t-t_q)/(6\cdot 365)})]$, con severidad
$\varsigma = \mathrm{clip}(0{,}72 + 0{,}3(1-g_{\min}) + 0{,}08\,G_2;\, 0{,}4;\, 0{,}98)$. La observación aplica
la ecuación~\eqref{eq:wcm} con el término de suelo modulado por la humedad estacional,
$\sigma^0_s\to\sigma^0_s\,10^{\Delta\sigma^0_{suelo}(t)/10}$ (ecuación~\ref{eq:dsigma}).

\subsection{Terremotos}
La deformación cosísmica $\vect u = (u_e, u_n, u_z)$ se calcula con la solución analítica de Okada
\cite{okada1985} para una falla rectangular en un semiespacio elástico ($\nu = 0{,}25$), con rumbo,
buzamiento, ángulo de deslizamiento, longitud $L_f$, ancho $W_f$, deslizamiento $s$ y profundidad del
centroide (la falla se mantiene al menos $0{,}2$~km bajo la superficie). La magnitud del escenario es
\begin{equation}
  M_w = \tfrac23\left(\log_{10}M_0 - 9{,}1\right),\qquad M_0 = \mu L_f W_f s,\quad \mu = 30~\mathrm{GPa}.
\end{equation}
La evolución temporal incluye el salto cosísmico y un término postsísmico logarítmico,
\begin{equation}
  f(t) = \begin{cases} 0 & t < t_s,\\ 1 + p_s\,\dfrac{\ln\!\left(1 + (t-t_s)/20\right)}{\ln 11} & t\ge t_s,\end{cases}
\end{equation}
con $p_s = 0{,}08$. La pila observada es $y_k = f(t_k)\,d_{LOS} + a_k - a_0 + n_k$, con
$n_k\sim\mathcal N(0,\sigma_d^2)$ y coherencia base $\gamma_0 = \mathrm{clip}(0{,}78 - 0{,}18\,G_5;\,0{,}3;\,0{,}95)$.
El océano (al oeste de una costa recta en el escenario de Pedernales) no tiene datos.

\subsection{Subsidencia}
La velocidad vertical es una suma de conos gaussianos de hundimiento con heterogeneidad multiplicativa,
\begin{equation}
  v(\vect x) = \left(1 + 0{,}2\,G_6\right)\sum_k f_r\,r_k\exp\!\left(-\frac{|\vect x - \vect c_k|^2}{2(f_R R_k)^2}\right),
\end{equation}
y el desplazamiento vertical acumulado es
\begin{equation}\label{eq:subsidencia}
  u_z(t) = v\,t_a + A_s\,\frac{|v|}{\max|v|}\,c(t),\qquad
  c(t) = \begin{cases}\tilde m(t) - \tilde m(t_0) & \text{con forzamiento},\\ \sin(2\pi t_a) & \text{sin forzamiento},\end{cases}
\end{equation}
donde $\tilde m$ es la anomalía normalizada de humedad del suelo (sección~\ref{sec:acople}). La
observación es $y_k = e_{up}\,u_z(t_k) + a_k - a_0 + n_k$.

\subsection{Glaciares}
El glaciar fluye por un canal sinuoso (amplitud 1200~m, longitud de onda 16~km) con coordenadas
curvilíneas $(s,\eta)$. El perfil transversal combina deslizamiento basal y deformación interna con la
ley de Glen de exponente $n=3$ \cite{glen1955}: $\Pi(\eta) = 0{,}5 + 0{,}5\,[1 - (|\eta|/W)^{n+1}]$. La
velocidad en la línea central crece hacia el frente $x_f(t)$,
\begin{equation}
  u_c(s) = u_i + (u_t-u_i)\exp\!\left(-\frac{\max(x_f - s, 0)}{L_a}\right),\quad
  x_f(t) = x_{f0} - r_f\,t_a - a_f\cos\frac{2\pi(t-t_p)}{365{,}25},
\end{equation}
y se multiplica por $[1 + \alpha(1-e^{-t/36})]\,[1 + 0{,}08\cos(2\pi(t-t_p)/365{,}25)]$. Cada componente
del desplazamiento recibe ruido gaussiano de desviación $\sigma_v$. Este modelo \emph{no} recibe forzamiento
estacional (sección~\ref{sec:acople}).

\subsection{Agricultura}
El área se divide en parcelas rectangulares separadas por caminos de un píxel. Cada parcela sigue una
rotación elegida al azar con probabilidades dadas; cada ciclo de cultivo con siembra $t_s$, cosecha
$t_c$, duración $D = t_c - t_s$ y biomasa máxima $B_{\max}$ crece de forma logística,
\begin{equation}
  B(t) = \frac{B_{\max}}{1 + \exp\!\left[-\dfrac{t - (t_s + 0{,}55D)}{D/10}\right]},\qquad t_s\le t<t_c,
\end{equation}
deja un rastrojo $0{,}8\,e^{-(t-t_c)/40}$ tras la cosecha y, en el arroz, una lámina de agua entre $t_s$ y
$t_c - 20$ días. La observación usa la ecuación~\eqref{eq:wcm} con parámetros de cultivo, un término
de suelo multiplicado por la humedad y por $0{,}03$ si la parcela está inundada, y un doble rebote
tallo--agua de $0{,}30\,B/B_{\max}$ (HH) y $0{,}006\,B/B_{\max}$ (HV). En la cuenca baja del Guayas
(Daule) la rotación dominante tiene dos ciclos de arroz: de invierno (siembra a mediados de enero,
cosecha en mayo) y de verano (mediados de junio a octubre).

%=====================================================================================================
\section{Detección de cambios}
%=====================================================================================================

\paragraph{Humedales.} Con $\overline{\mathrm{HH}}$ y $\overline{\mathrm{HV}}$ promediados $3\times3$ (dB),
un píxel está inundado si $\overline{\mathrm{HH}}<-16$ (agua abierta), si $\overline{\mathrm{HH}}>-5{,}5$ y
$\overline{\mathrm{HH}}-\overline{\mathrm{HV}}>6$ (doble rebote en bosque inundado) o si
$-9{,}5<\overline{\mathrm{HH}}\le-5{,}5$ y $\overline{\mathrm{HH}}-\overline{\mathrm{HV}}>8{,}5$ (herbácea
inundada). La exactitud se mide píxel a píxel contra la verdad del modelo.

\paragraph{Incendios.} La razón logarítmica
$LR = 10\log_{10}\left(\langle\sigma^0_{HV}\rangle_{1}/\langle\sigma^0_{HV}\rangle_{0}\right)$ marca como
quemado el píxel con $LR<-2{,}5$~dB y $\langle\sigma^0_{HV}\rangle_0>-18$~dB (severidad alta si
$LR<-4{,}5$~dB). Las emisiones se estiman como $\mathrm{CO_2}\approx\Delta B\cdot 0{,}47\cdot 44/12$.

\paragraph{Terremotos.} El interferograma entre $i_0$ e $i_1$ usa la coherencia
$\gamma = \gamma_0\,(0{,}55 + 0{,}45\,e^{-\Delta t/90})$, reducida cerca de la traza de la falla por
$1 - 0{,}9\,e^{-(D_f/1{,}5\Delta x)^2}$ si el sismo ocurre en el intervalo. Se enmascaran los píxeles con
$\gamma<0{,}25$, como hacen las componentes conexas del producto GUNW.

\paragraph{Subsidencia.} En cada píxel se ajusta por mínimos cuadrados
$d(t) = a + v\,(t-\bar t) + A\sin 2\pi t + B\cos 2\pi t$ (los términos armónicos solo si la serie supera
$0{,}8$ años) y se convierte $v_{vert} = v_{LOS}/e_{up}$ con un multilook $3\times3$.

\paragraph{Glaciares.} La velocidad es $|\Delta\vect r|/\Delta t$ entre dos fechas (producto GOFF); el
frente se localiza sobre la línea de flujo central como el último punto con velocidad superior al umbral
de hielo, tolerando huecos de hasta cuatro muestras.

\paragraph{Agricultura.} Por parcela se calcula la media lineal de HV; la mayor caída entre fechas
consecutivas menor que $-3$~dB indica cosecha, o inundación para trasplante si el nivel posterior es
menor que $-32$~dB; un aumento mayor que $3$~dB indica crecimiento.

%=====================================================================================================
\section{Estaciones y temporadas en Ecuador}\label{sec:estaciones}
%=====================================================================================================

Ecuador no tiene cuatro estaciones térmicas: la temperatura depende sobre todo de la altitud y el ciclo
anual dominante es el de la lluvia, con calendarios distintos por región \cite{pourrut1995,vuille2000}.
La Costa tiene una época lluviosa (llamada popularmente \emph{invierno}) entre diciembre--enero y
abril--mayo y una época seca (\emph{verano}) el resto del año; los valles interandinos suelen tener un
régimen bimodal; la Amazonía es húmeda todo el año; y en Galápagos la corriente de Humboldt alterna una
época cálida y lluviosa con otra fría y seca de garúa. El módulo \codigo{modelo/estaciones.py} cuantifica estas
diferencias en tres niveles: astronómico, climatológico e hidrológico.

\subsection{Estación astronómica}
La declinación solar del día juliano $J$ se calcula con la serie de Spencer \cite{spencer1971}
(error menor que $0{,}035^\circ$), con $\Gamma = 2\pi(J-1)/365$:
\begin{equation}
\begin{aligned}
  \delta = {}&0{,}006918 - 0{,}399912\cos\Gamma + 0{,}070257\sin\Gamma - 0{,}006758\cos2\Gamma\\
  &+ 0{,}000907\sin2\Gamma - 0{,}002697\cos3\Gamma + 0{,}00148\sin3\Gamma .
\end{aligned}
\end{equation}
Los equinoccios son los ceros de $\delta(J)$ y los solsticios sus extremos (tabla~\ref{tab:hitos}); al
ignorar la ecuación del tiempo y el año bisiesto, las fechas resultantes pueden diferir en un día de las
efemérides astronómicas, precisión suficiente para nombrar la estación. La
estación astronómica se nombra según el hemisferio del punto (en Ecuador continental, casi siempre el
austral). El ángulo horario de la puesta, la duración del día, la elevación del Sol al mediodía y la
radiación extraterrestre son \cite{allen1998}
\begin{gather}
  \omega_s = \arccos(-\tan\varphi\tan\delta),\qquad N = \frac{24}{\pi}\,\omega_s,\qquad
  \epsilon_{12} = 90^\circ - |\varphi-\delta|,\\
  R_a = \frac{24\cdot60}{\pi}\,G_{sc}\,d_r\left[\omega_s\sin\varphi\sin\delta + \cos\varphi\cos\delta\sin\omega_s\right],
  \qquad d_r = 1 + 0{,}033\cos\frac{2\pi J}{365},
\end{gather}
con $G_{sc} = 0{,}0820$~MJ\,m$^{-2}$\,min$^{-1}$. En la zona tropical el Sol pasa por el cenit dos veces al
año, en las fechas en que $\delta(J) = \varphi$; cerca del ecuador ambas caen próximas a los equinoccios y
$N$ apenas varía (tabla~\ref{tab:hitos}). Por eso, en Ecuador, la estación astronómica casi no modula la
energía disponible y la temporada relevante es la de lluvias.

\TablaHitos

\subsection{Climatología NASA POWER}
Se usa la climatología mensual 2001--2025 del proyecto NASA POWER \cite{power2025} (API de climatología
v2.10), que deriva las variables meteorológicas del reanálisis MERRA-2 \cite{gelaro2017} en su malla
nativa de $0{,}5^\circ$ en latitud por $0{,}625^\circ$ en longitud. Se descargan 230 nodos que cubren el
Ecuador continental ($-5{,}5^\circ$ a $2^\circ$, $-81{,}25^\circ$ a $-75^\circ$) y Galápagos, con las
variables: precipitación total corregida por sesgo $P$ (\codigo{PRECTOTCORR}, mm\,día$^{-1}$), temperatura
media a 2~m $\bar T$ (\codigo{T2M}), rango térmico diario medio $\Delta T$ (\codigo{T2M\_RANGE}) y agua
precipitable $PW$ (\codigo{PW}, cm). El agua precipitable se distribuye en una malla más gruesa y se
asigna a cada nodo desde el nodo más cercano. En un punto $(\varphi,\lambda)$ cada variable mensual se
interpola bilinealmente:
\begin{equation}
  q(\varphi,\lambda) = \sum_{a,b\in\{0,1\}} \left[(1-u)^{1-a}u^{a}\right]\left[(1-v)^{1-b}v^{b}\right] q_{i+a,\,j+b},
\end{equation}
con $u,v\in[0,1)$ las coordenadas fraccionarias dentro de la celda. Si falta alguna de las cuatro
esquinas se usa el nodo más cercano, siempre que esté a menos de $1^\circ$.

\subsection{Temporada climática}
Con $P_m$ la precipitación del mes $m$ (mm) y $\bar P$ la media mensual, cada mes se clasifica como
\begin{equation}
  \text{temporada}(m) = \begin{cases}
    \text{seca} & P_m < 60~\text{mm},\\
    \text{lluviosa} & P_m \ge \bar P,\\
    \text{de transición} & \text{en otro caso},
  \end{cases}
\end{equation}
donde el umbral de 60~mm es el criterio de mes seco de la clasificación climática tropical de
Köppen--Geiger \cite{peel2007}. El régimen es unimodal o bimodal según el número de rachas circulares
de meses con $P_m\ge\bar P$. En Galápagos la temporada se nombra además por la temperatura: época cálida
si $\bar T_m$ supera la media anual y época fría de garúa en otro caso.

\TablaClimaSitios

\FiguraClimaSitios

\subsection{Evapotranspiración de referencia}
Como solo se dispone de temperatura, la evapotranspiración de referencia se estima con la ecuación de
Hargreaves--Samani \cite{hargreaves1985}, recomendada por FAO-56 para ese caso \cite{allen1998}:
\begin{equation}
  ET_0 = 0{,}0023\,(\bar T + 17{,}8)\,\sqrt{\Delta T}\;0{,}408\,R_a\quad[\text{mm\,día}^{-1}],
\end{equation}
con $R_a$ evaluada en el día central de cada mes.

\subsection{Balance hídrico diario del suelo}
Los valores mensuales de $P$ y $ET_0$ se interpolan linealmente y de forma periódica entre los días
centrales de cada mes para obtener series diarias $P_t$, $ET_{0,t}$. El agua del suelo $W$ (mm) sigue un
modelo de cubeta de Thornthwaite--Mather \cite{thornthwaite1955}, con evapotranspiración real
proporcional al agua disponible:
\begin{equation}
\begin{aligned}
  ET_{a,t} &= ET_{0,t}\,\frac{W_t}{W_{\max}},\qquad W^\ast = W_t + P_t - ET_{a,t},\\
  Q_t &= \max(W^\ast - W_{\max},\,0),\qquad W_{t+1} = \min\!\left(\max(W^\ast,0),\, W_{\max}\right),
\end{aligned}
\end{equation}
con $W_{\max} = 150$~mm, el agua disponible total de un suelo franco con un metro de raíces
\cite{allen1998}. Partiendo de $W_0 = W_{\max}/2$ se repite el año climatológico seis veces; el último año
es el ciclo periódico estable. La humedad relativa del suelo es $s_t = W_t/W_{\max}\in[0,1]$ y $Q_t$ es la
escorrentía por saturación.

\subsection{Almacenamiento de un humedal}
El humedal se representa como un embalse lineal alimentado por el excedente del suelo:
\begin{equation}\label{eq:embalse}
  \frac{dS}{dt} = Q - \frac{S}{\tau},\qquad S_{t+1} = S_t + Q_t - \frac{S_t}{\tau}\quad(\Delta t = 1~\text{día}),
\end{equation}
con $\tau = 30$~días, integrado con el mismo periodo de arranque. El esquema de Euler explícito es estable
porque $\Delta t/\tau \ll 1$.

\subsection{Acoplamiento con el modelo radar}\label{sec:acople}
Para una serie $x_t$ se define la anomalía normalizada
\begin{equation}
  \tilde x_t = \frac{x_t - \bar x}{\max_t|x_t - \bar x|}\in[-1,1]
\end{equation}
($\tilde x\equiv 0$ si la serie no varía). Todas las series se evalúan en cualquier instante
convirtiendo $t$ a día juliano fraccionario e interpolando periódicamente; así el forzamiento queda
ligado al calendario y no a la fecha de referencia del escenario. Los acoples son:
\begin{enumerate}
  \item \textbf{Retrodispersión del suelo.} La humedad volumétrica en exceso sobre la media es
        $\Delta\theta = \theta_u\,(s - \bar s)$, con $\theta_u = 0{,}15$~m$^3$\,m$^{-3}$ (el agua útil de
        150~mm en un metro). El cambio de retrodispersión del suelo es
        \begin{equation}\label{eq:dsigma}
          \Delta\sigma^0_{suelo} = \kappa\cdot 100\,\Delta\theta,\qquad \kappa = 0{,}25~\mathrm{dB}/\%,
        \end{equation}
        dentro del intervalo de sensibilidades de $\sigma^0$ a la humedad volumétrica en suelo desnudo
        descrito por Ulaby et al.\ \cite{ulaby1978,ulaby2014}. Se aplica al término de suelo del modelo de
        nube de agua en incendios y agricultura (en esta última, más una heterogeneidad de riego
        $U(-1,1)$~dB por parcela).
  \item \textbf{Pulso de inundación.} En la ecuación~\eqref{eq:nivel}, $p(t) = \tilde S(t)$.
  \item \textbf{Retardo troposférico.} El retardo cenital húmedo es proporcional al agua precipitable
        ($ZWD\approx PW/\Pi$, $\Pi\approx 0{,}15$) \cite{bevis1992}; en primera aproximación se escala la
        desviación de la pantalla atmosférica con
        \begin{equation}\label{eq:fpw}
          F_{pw}(t) = PW(t)/\overline{PW}.
        \end{equation}
  \item \textbf{Deformación estacional.} En la ecuación~\eqref{eq:subsidencia}, $c(t) = \tilde m(t) -
        \tilde m(t_0)$ con $\tilde m = \tilde s$: los rellenos arcillosos se hinchan y el acuífero se
        recarga en la época lluviosa \cite{galloway2011}.
  \item \textbf{Calendario agrícola.} Las fechas de las rotaciones se definen respecto de una fecha fija
        del año civil (15 de mayo en Daule) y se replican el año anterior y el siguiente:
        $t_s = t_s^{(0)} + \delta_j + (J_{ref} - J_{evento}) + 365\,a$, $a\in\{-1,0,1\}$, con $\delta_j$ el
        desfase aleatorio de cada parcela.
  \item \textbf{Glaciares.} Sin forzamiento: en los glaciares ecuatoriales la ablación ocurre todo el
        año y la estacionalidad del flujo es débil e incierta \cite{francou2004,rabatel2013}.
\end{enumerate}
La figura~\ref{fig:forzamiento} muestra el ciclo resultante en La Segua: el suelo se satura en febrero,
el humedal alcanza su máximo a finales de marzo y se vacía durante la época seca, y la retrodispersión
del suelo oscila entre unos $-1{,}6$ y $+1{,}8$~dB respecto de su media.

\FiguraForzamiento

%=====================================================================================================
\section{Verificación}
%=====================================================================================================

La tabla~\ref{tab:verificacion} compara los seis escenarios de Ecuador sin y con forzamiento. En La
Segua el máximo de inundación pasa del 31 de marzo (fase fija heredada) al 19 de marzo, dentro de la
época lluviosa costera. En Loja, la amplitud estacional de HV en el bosque no quemado es de solo
$\approx0{,}1$~dB en ambos casos: con 120~Mg\,ha$^{-1}$ la transmisividad $\tau^2\approx 0{,}15$ deja que el
dosel domine la señal, de modo que la humedad del suelo apenas afecta a la detección. En Pedernales la
desviación atmosférica varía entre adquisiciones en un $\pm6$~\% siguiendo el agua precipitable; en
Guayaquil el error de la velocidad estimada cambia ligeramente al sustituir la sinusoide arbitraria por
el ciclo hídrico. Los glaciares no cambian por diseño.

\TablaVerificacion

La tabla~\ref{tab:invariancia} es la prueba decisiva de la corrección: al desplazar la fecha de referencia
del escenario de La Segua 90 y 180 días, el máximo del pulso de inundación se desplaza igual con la fase
fija (de finales de marzo a finales de junio y de septiembre, en plena época seca), mientras que con el
forzamiento permanece en marzo. Del mismo modo, las cosechas de arroz de Daule siguen cayendo en
abril--mayo y septiembre--octubre aunque la fecha de referencia se mueva cuatro meses.

\TablaInvariancia

%=====================================================================================================
\section{Limitaciones e incertidumbre}
%=====================================================================================================

\begin{itemize}
  \item \textbf{Resolución.} La malla MERRA-2 de $\sim$55~km no resuelve los valles interandinos ni los
        gradientes de las estribaciones; por ejemplo, en Quito el régimen bimodal observado
        \cite{pourrut1995} aparece como unimodal en la malla. La elevación de la celda ($z_c$ en la tabla~\ref{tab:clima-sitios}) puede diferir en
        miles de metros de la del sitio.
  \item \textbf{Climatología, no año real.} Se usa el ciclo medio 2001--2025: no se representan la
        variabilidad interanual ni los eventos de El Niño y La Niña, que en la Costa pueden multiplicar la
        lluvia de la época húmeda \cite{vuille2000}.
  \item \textbf{Evapotranspiración.} La ecuación de Hargreaves--Samani tiende a sobrestimar $ET_0$ en
        climas húmedos \cite{droogers2002}, lo que acelera el secado del suelo en la Costa.
  \item \textbf{Balance hídrico.} Una sola capa con $W_{\max}$ uniforme, sin enrutamiento de
        escorrentía ni aporte de la cuenca al humedal más allá del excedente local.
  \item \textbf{Sensibilidad radar.} $\kappa$ depende de la incidencia, la rugosidad y la textura; el
        valor adoptado es representativo, no calibrado para cada sitio.
  \item \textbf{Atmósfera.} El escalado por $PW$ ignora que la turbulencia convectiva de la época
        lluviosa puede aumentar la varianza más que la media del retardo.
  \item \textbf{Región natural.} Es la de la provincia (división administrativa), no un piso climático.
\end{itemize}

%=====================================================================================================
\section{Implementación y reproducibilidad}
%=====================================================================================================

\begin{center}\small
\begin{tabular}{@{}lp{9.6cm}@{}}\toprule Ruta & Contenido \\ \midrule
\codigo{app.py} & Punto de entrada: interfaz Streamlit (comentada línea a línea) con mapa restringido a
  Ecuador, reproductor, antes/después, pestaña \emph{Estación y temporada}, catálogo ASF y explorador. \\
\codigo{modelo/nisar.py} & Geometría, modelos directos, simulación, detección y lectura de HDF5 GCOV/GUNW. \\
\codigo{modelo/estaciones.py} & Regiones naturales, astronomía solar, climatología, temporadas, balance hídrico
  y forzamiento estacional (\codigo{Forzamiento}). \\
\codigo{modelo/clima.py} & Partición espacial reproducible para validar clasificadores (otro informe). \\
\codigo{interfaz/tema.py} & Tema cyberpunk: paleta neón validada, rampas de color, CSS y estilo de gráficos. \\
\codigo{scripts/preparar\_datos\_ecuador.py} & Descarga y prepara \codigo{data/ecuador\_provincias.geojson} y
  \codigo{data/clima\_ecuador\_power.csv}. \\
\codigo{docs/generar\_datos\_latex.py} & Calcula las tablas y figuras de los informes. \\
\codigo{docs/estilo\_cyberpunk.sty} & Maqueta común de los informes (opción \codigo{impresion}: fondo blanco). \\
\bottomrule
\end{tabular}
\end{center}

\noindent Secuencia para reproducir los resultados (desde la raíz del proyecto):
\begin{cpcaja}[fontupper=\small]
\begin{verbatim}
python scripts/preparar_datos_ecuador.py  # datos (NASA POWER, geoBoundaries)
python docs/generar_datos_latex.py        # tablas y figuras de los informes
cd docs
pdflatex modelo_nisar_ecuador.tex         # dos pasadas: índice y referencias
pdflatex modelo_nisar_ecuador.tex
cd ..
python app.py                             # aplicación en http://localhost:8501
\end{verbatim}
\end{cpcaja}
Las simulaciones son deterministas para una semilla dada (generador \codigo{default\_rng} de NumPy).

%=====================================================================================================
\begin{thebibliography}{99}\small

\bibitem{allen1998} Allen, R.~G., Pereira, L.~S., Raes, D. y Smith, M. (1998). \emph{Crop
evapotranspiration: Guidelines for computing crop water requirements}. FAO Irrigation and Drainage Paper
56. FAO, Roma.

\bibitem{attema1978} Attema, E.~P.~W. y Ulaby, F.~T. (1978). Vegetation modeled as a water cloud.
\emph{Radio Science}, 13(2), 357--364.

\bibitem{bevis1992} Bevis, M., Businger, S., Herring, T.~A., Rocken, C., Anthes, R.~A. y Ware, R.~H.
(1992). GPS meteorology: Remote sensing of atmospheric water vapor using the Global Positioning System.
\emph{Journal of Geophysical Research}, 97(D14), 15787--15801.

\bibitem{droogers2002} Droogers, P. y Allen, R.~G. (2002). Estimating reference evapotranspiration under
inaccurate data conditions. \emph{Irrigation and Drainage Systems}, 16, 33--45.

\bibitem{francou2004} Francou, B., Vuille, M., Favier, V. y Cáceres, B. (2004). New evidence for an ENSO
impact on low-latitude glaciers: Antizana 15, Andes of Ecuador, 0$^\circ$28'S. \emph{Journal of
Geophysical Research}, 109, D18106.

\bibitem{galloway2011} Galloway, D.~L. y Burbey, T.~J. (2011). Review: Regional land subsidence
accompanying groundwater extraction. \emph{Hydrogeology Journal}, 19, 1459--1486.

\bibitem{gelaro2017} Gelaro, R. et al. (2017). The Modern-Era Retrospective Analysis for Research and
Applications, Version 2 (MERRA-2). \emph{Journal of Climate}, 30(14), 5419--5454.

\bibitem{glen1955} Glen, J.~W. (1955). The creep of polycrystalline ice. \emph{Proceedings of the Royal
Society of London A}, 228, 519--538.

\bibitem{goodman1976} Goodman, J.~W. (1976). Some fundamental properties of speckle. \emph{Journal of the
Optical Society of America}, 66(11), 1145--1150.

\bibitem{hanssen2001} Hanssen, R.~F. (2001). \emph{Radar Interferometry: Data Interpretation and Error
Analysis}. Kluwer Academic Publishers, Dordrecht.

\bibitem{hargreaves1985} Hargreaves, G.~H. y Samani, Z.~A. (1985). Reference crop evapotranspiration from
temperature. \emph{Applied Engineering in Agriculture}, 1(2), 96--99.

\bibitem{power2025} NASA Langley Research Center (2025). \emph{POWER Project: Prediction Of Worldwide
Energy Resources}, API de climatología v2.10. \url{https://power.larc.nasa.gov} (consulta: 1 de noviembre
de 2026).

\bibitem{nisar2018} NISAR (2018). \emph{NASA-ISRO SAR (NISAR) Mission Science Users' Handbook}. NASA Jet
Propulsion Laboratory, Pasadena.

\bibitem{okada1985} Okada, Y. (1985). Surface deformation due to shear and tensile faults in a
half-space. \emph{Bulletin of the Seismological Society of America}, 75(4), 1135--1154.

\bibitem{peel2007} Peel, M.~C., Finlayson, B.~L. y McMahon, T.~A. (2007). Updated world map of the
Köppen--Geiger climate classification. \emph{Hydrology and Earth System Sciences}, 11, 1633--1644.

\bibitem{pourrut1995} Pourrut, P. (ed.) (1995). \emph{El agua en el Ecuador: clima, precipitaciones,
escorrentía}. ORSTOM y Corporación Editora Nacional, Quito.

\bibitem{rabatel2013} Rabatel, A. et al. (2013). Current state of glaciers in the tropical Andes: a
multi-century perspective on glacier evolution and climate change. \emph{The Cryosphere}, 7, 81--102.

\bibitem{rodriguez1992} Rodríguez, E. y Martin, J.~M. (1992). Theory and design of interferometric
synthetic aperture radars. \emph{IEE Proceedings F}, 139(2), 147--159.

\bibitem{runfola2020} Runfola, D. et al. (2020). geoBoundaries: A global database of political
administrative boundaries. \emph{PLoS ONE}, 15(4), e0231866.

\bibitem{spencer1971} Spencer, J.~W. (1971). Fourier series representation of the position of the Sun.
\emph{Search}, 2(5), 172.

\bibitem{thornthwaite1955} Thornthwaite, C.~W. y Mather, J.~R. (1955). The water balance.
\emph{Publications in Climatology}, 8(1). Drexel Institute of Technology, Centerton.

\bibitem{ulaby1978} Ulaby, F.~T., Batlivala, P.~P. y Dobson, M.~C. (1978). Microwave backscatter
dependence on surface roughness, soil moisture, and soil texture: Part~I---Bare soil. \emph{IEEE
Transactions on Geoscience Electronics}, GE-16(4), 286--295.

\bibitem{ulaby2014} Ulaby, F.~T. y Long, D.~G. (2014). \emph{Microwave Radar and Radiometric Remote
Sensing}. University of Michigan Press, Ann Arbor.

\bibitem{vuille2000} Vuille, M., Bradley, R.~S. y Keimig, F. (2000). Climate variability in the Andes of
Ecuador and its relation to tropical Pacific and Atlantic sea surface temperature anomalies.
\emph{Journal of Climate}, 13, 2520--2535.

\bibitem{williams1998} Williams, S., Bock, Y. y Fang, P. (1998). Integrated satellite interferometry:
Tropospheric noise, GPS estimates and implications for interferometric synthetic aperture radar
products. \emph{Journal of Geophysical Research}, 103(B11), 27051--27067.

\end{thebibliography}
\end{document}
